\documentclass[11pt]{article}
\usepackage[margin=0.85in]{geometry}
\usepackage{lmodern}
\usepackage{microtype}
\usepackage{enumitem}
\usepackage{listings}
\usepackage[T1]{fontenc}

\lstset{
  basicstyle=\ttfamily\small,
  columns=fullflexible,
  keepspaces=true,
  showstringspaces=false,
  breaklines=true,
  frame=single,
  xleftmargin=0.4em,
  xrightmargin=0.4em,
  aboveskip=0.55em,
  belowskip=0.55em
}

\pagestyle{empty}
\setlength{\parskip}{0.35em}
\setlength{\parindent}{0pt}
\setlist[enumerate]{leftmargin=1.4em, itemsep=0.65em, topsep=0.3em}
\setlist[itemize]{leftmargin=1.3em, itemsep=0.15em, topsep=0.2em}

\begin{document}

\begin{center}
{\LARGE\bfseries Quiz 3 --- Sample 2}\\[0.35em]
DATA 412/612
\end{center}

\vspace{0.25em}

Weeks 1--10. The new material is stringr, factors, and lubridate.

Write your answers in an R script. You may use R and RStudio.
You may use your notes and the course website
(\texttt{https://amousavi1.github.io/612/data-612.html}).
You may not use Google, any other website, or an LLM.
Load the tidyverse.

\begin{enumerate}
\item In \texttt{people}, turn the missing value into the text
\texttt{missing}, put every string in title case, and paste the
three strings into one string separated by \texttt{"; "}.

\begin{lstlisting}
people <- c("ada lovelace", "grace hopper", NA)
\end{lstlisting}

\item \texttt{party} is already a factor. Merge \texttt{Dem} and
\texttt{Rep} into one level named \texttt{Major}. Then put the
most common level first. Report the levels.

\begin{lstlisting}
party <- factor(c(
  "Dem", "Dem", "Rep", "Rep", "Rep", "Ind", "Other", "Other"
))
\end{lstlisting}

\item Start from \texttt{ymd("2016-01-01")}. The year 2016 is a
leap year. What date does adding \texttt{years(1)} land on?
What date does adding \texttt{ddays(365)} land on?
Which of those two functions is a period, and which is a duration?
Also give \texttt{floor\_date(ymd("2026-11-11"), unit = "month")}.

\item In \texttt{events}, the date strings are day, then month,
then year. Parse them. Keep rows on or after 1 February 2010.
Put \texttt{label} in upper case, then set the level order to
\texttt{TALK} followed by \texttt{LAB-A}. Report the levels, and
the labels in row order.

\begin{lstlisting}
events <- tribble(
  ~label, ~opened,
  "Lab-A", "01-02-2010",
  "Lab-B", "15-01-2010",
  "Talk",  "10-03-2010"
)
\end{lstlisting}
\end{enumerate}

\end{document}
